\documentclass[10pt,a4paper]{amsart}
\usepackage[margin=19mm]{geometry}
\usepackage{amsmath,amssymb,mathtools}
\usepackage[scr=rsfs,cal=euler]{mathalfa}
\usepackage{xfrac}
\usepackage[unicode,hidelinks,bookmarks=false]{hyperref}
\newcommand{\ie}{i.e.\ }
\newcommand{\cf}{cf.\ }
\newcommand{\resp}{respectively\ }
\newcommand{\Subsection}{\subsection}
\providecommand{\nobreakdash}{\nobreak\hbox{-}}
\setlength{\emergencystretch}{2em}
\multlinegap0pt
\usepackage{amssymb}
\usepackage[shortlabels]{enumitem}
\newtheorem{theorem}{Theorem}
\newtheorem{proposition}{Proposition}
\title{The converse to Borsuk’s result on fans fails}
\author{Benjamin Vejnar}
\date{Source: arXiv:2604.12517v1 (CC BY 4.0)}
\begin{document}
\maketitle

\noindent\emph{Source and licence.} The converse to Borsuk’s result on fans fails. Version arXiv:2604.12517v1 (CC BY 4.0), distributed by arXiv under the Creative Commons Attribution 4.0 International licence, http://creativecommons.org/licenses/by/4.0/, as recorded in the arXiv licence metadata for this exact version and shown on the abstract page https://arxiv.org/abs/2604.12517v1. Full licence notice: \texttt{licenses/arxiv-2604.12517-CC-BY-4.0.txt}.\par\smallskip

\noindent\textit{Editorial scope.} Counted unit: the article's proposition proposition (source0.tex lines 490--493). The article's complete original proof of it is delivered with it (source0.tex lines 495--503, one \texttt{proof} environment of 9 lines). No mathematics is written, repaired or re-derived: the statement and the proof are the article's own lines, in the article's own notation, and no retained line is substituted or re-flowed.  Retained uncounted as the source-local context the proof needs: lines 476--478.  Nothing was omitted from any retained span, so the package carries no omission record.  No retained line contains a citation, so the wrapper carries no bibliography and no bibliography entry.  No retained line contains a figure environment, an image-inclusion command or drawing code, so no image is delivered and the package declares no asset dependency.  Editorial formatting only, in addition: an AMS \texttt{amsart} preamble that restates the article's own declarations and macros used by the retained lines, namely the environments \texttt{theorem}, \texttt{proposition} on separate counters, together with the standard packages those lines need.  The retained environments print in this document as Theorem 1, Proposition 1 in order of appearance, whereas the article prints the same items with the numbering of its own full document.  The complete native source of the article as distributed in the arXiv e-print for this exact version is bundled unchanged as \texttt{sources/198406/source0.tex}.
\begin{theorem}\label{TheoremZeroOneLaw}
For every $h\in\mathbb H$ the set $\mathbb L_h$ is either meager or comeager.
\end{theorem}

\begin{proposition}
For the generic Cantor set homeomorphism $h\in\mathbb H$ the set $\mathbb L_h$ is meager, i.e. $h$ can be extended only to meager many Lelek fences $L_\varphi$, for $\varphi\in\mathbb L$.
(Consequently, for every $\varphi\in\mathbb L$, the set of those $h\in\mathbb H$ which can be linearly extended to $L_\varphi$ is meager).
\end{proposition}

\begin{proof}
Suppose that $h\in\mathbb H$ is the generic homeomorphism of the Cantor set and suppose for contradiction that $\mathbb L_h$ is not meager. Then by Theorem \ref{TheoremZeroOneLaw} it follows that $\mathbb L_h$ is comeager.
Consider the set 
\[S=\{(h,\varphi)\in\mathbb H\times\mathbb L: \varphi\in\mathbb L_h\}.\]
The set $S$ is Baire measurable (since analytic)
and its vertical slices $S_g=\mathbb L_g$ are comeager for every $g$ conjugate to $h$. Hence comeager many vertical slices are comeager and thus by the Kuratowski-Ulam theorem the set $S$ is comeager.
Consequently, by applying the Kuratowski-Ulam theorem in the horizontal direction, there exists $\psi\in\mathbb L$ for which the set $G:=\{g\in\mathbb H: (g, \psi)\in S\}$ is comeager.
Thus $G$ is a comeager subgroup of the Polish group $\mathbb H$ and thus $G=\mathbb H$. This is a contradiction.
\end{proof}

\end{document}
